%%=====================================================================
%%  preamble.tex — SSRN preprint build (plain article class).
%%  Sections are the SUBMITTED JFEC text (JFEC_OUP/sections, 2026-09-02),
%%  copied verbatim. Only this wrapper differs from the OUP build.
%%=====================================================================
\usepackage[margin=1in]{geometry}
\usepackage[T1]{fontenc}
\usepackage[utf8]{inputenc}
\usepackage{microtype}
\usepackage{lmodern}
\usepackage{setspace}
\usepackage{amsmath,amssymb,amsthm}
\usepackage{mathtools}
\usepackage{bm}
\usepackage{booktabs}
\usepackage{graphicx}
\graphicspath{{figures/}}
\usepackage{array}
\usepackage{multirow}
\usepackage{rotating}
\usepackage{float}
\usepackage{threeparttable}
\usepackage{xcolor}
\usepackage{caption}
\captionsetup{font=small,labelfont=bf}
\usepackage{appendix}  % no [title]: the OUP build prints "A. Appendix: ..." too
\usepackage{tikz}
\usetikzlibrary{positioning,calc,fit,backgrounds,arrows.meta,shapes.geometric,decorations.pathreplacing,patterns}
\usepackage[round,authoryear]{natbib}
\usepackage[hidelinks]{hyperref}
\hypersetup{pdftitle={Quantifying the Contributions of Clustering to Statistical Arbitrage},pdfauthor={Lunji Zhu, Yixuan He, Mihai Cucuringu}}
% float placement (the article class strands floats otherwise)
\setcounter{topnumber}{4}\setcounter{bottomnumber}{2}\setcounter{totalnumber}{5}
\renewcommand{\topfraction}{0.92}\renewcommand{\bottomfraction}{0.75}
\renewcommand{\textfraction}{0.06}\renewcommand{\floatpagefraction}{0.70}
% --- cross-reference helpers (cleveref is not installed here) --------
\newcommand{\secref}[1]{Section~\ref{#1}}
\newcommand{\secsref}[2]{Sections~\ref{#1} and~\ref{#2}}
\newcommand{\tabref}[1]{Table~\ref{#1}}
\newcommand{\figref}[1]{Figure~\ref{#1}}
\newcommand{\eqnref}[1]{Equation~\eqref{#1}}
\newcommand{\appref}[1]{Appendix~\ref{#1}}
\newcommand{\panelref}[2]{Panel~#1 of \tabref{#2}}


%%---------------------------------------------------------------------
%%  NOTATION — identical to the article-class build; the sections use it.
%%---------------------------------------------------------------------
\newcommand{\N}{N}
\newcommand{\Tw}{T}
\newcommand{\K}{K}
\newcommand{\Rmat}{\mathbf{R}}
\newcommand{\Cmat}{\mathbf{C}}
\newcommand{\Wmat}{\mathbf{W}}
\newcommand{\Dmat}{\mathbf{D}}
\newcommand{\Proj}{\mathbf{P}}
\newcommand{\Imat}{\mathbf{I}}
\newcommand{\rt}{\tilde{\mathbf{r}}}
\newcommand{\Zmat}{\mathbf{Z}}
\newcommand{\Lam}{\bm{\Lambda}}
\newcommand{\ones}[1]{\bm{1}_{#1}}
\newcommand{\Pgics}{\mathbf{P}_{\mathrm{GICS}}}
\newcommand{\SR}{\mathrm{SR}}
\newcommand{\Exp}{\mathbb{E}}
\newcommand{\Prob}{\mathbb{P}}
\newcommand{\Var}{\mathrm{Var}}
\newcommand{\Cov}{\mathrm{Cov}}
\newcommand{\real}{\mathbb{R}}
\newcommand{\tr}{\operatorname{tr}}
\newcommand{\rank}{\operatorname{rank}}
\DeclareMathOperator*{\argmin}{arg\,min}
\DeclareMathOperator*{\argmax}{arg\,max}
\newcommand{\Part}{\mathcal{P}}
\newcommand{\PartOne}{\Part_{\mathrm{one}}}
\newcommand{\PartRnd}{\Part_{\mathrm{rand}}}
\newcommand{\alphaFF}{\alpha_{\mathrm{FF5+MOM}}}
\newcommand{\ARI}{\mathrm{ARI}}
\newcommand{\nogroup}{\emph{no grouping}}
\newcommand{\randgroup}{\emph{random grouping}}
